\textbf{Phần 1. Trắc nghiệm lựa chọn}
\setcounterpageref{ex}{0}

\begin{ex} %[2H11N57]
Để đo chiều dài của phòng học, người ta nên dùng đơn vị nào?
\choice
{Mi-li-mét}
{Ki-lô-mét}
{Đề-xi-mét}
{\True Mét}
\loigiai{
Đo phòng học dùng đơn vị mét.\\
\textbf{Đáp án của câu hỏi là D}
}
\end{ex}

\begin{ex} %[2H11N57]
Để đo khoảng cách giữa hai thành phố, người ta dùng đơn vị
\choice
{\True Ki-lô-mét}
{Mét}
{Xăng-ti-mét}
{Đề-xi-mét}
\loigiai{
Đo khoảng cách giữa 2 thành phố dùng ki-lô-mét.\\
\textbf{Đáp án của câu hỏi là A}
}
\end{ex}

\begin{ex} %[2H11H57]
Bảng lớp dài khoảng bao nhiêu?
\choice
{$3\text{ cm}$}
{$3\text{ km}$}
{$30\text{ m}$}
{\True $3\text{ m}$}
\loigiai{
Bảng lớp dài khoảng $3\text{ m}$.\\
\textbf{Đáp án của câu hỏi là D}
}
\end{ex}

\begin{ex} %[2H11H57]
Bút chì của em dài khoảng bao nhiêu?
\choice
{$15\text{ m}$}
{$15\text{ dm}$}
{\True $15\text{ cm}$}
{$15\text{ km}$}
\loigiai{
Bút chì dài khoảng $15\text{ cm}$.\\
\textbf{Đáp án của câu hỏi là C}
}
\end{ex}

\begin{ex} %[2H11H57]
Cột cờ trường học cao khoảng
\choice
{$8\text{ cm}$}
{$8\text{ km}$}
{$80\text{ m}$}
{\True $8\text{ m}$}
\loigiai{
Cột cờ cao khoảng $8\text{ m}$.\\
\textbf{Đáp án của câu hỏi là D}
}
\end{ex}

\begin{ex} %[2H11H57]
Một gang tay của học sinh lớp 2 dài khoảng
\choice
{$15\text{ m}$}
{\True $15\text{ cm}$ (hoặc hơn 1 dm)}
{$15\text{ km}$}
{$1\text{ cm}$}
\loigiai{
Gang tay dài khoảng $15\text{ cm}$.\\
\textbf{Đáp án của câu hỏi là B}
}
\end{ex}

\begin{ex} %[2H11V57]
Đoạn thẳng $AB$ dài $1\text{ m}$, đoạn thẳng $CD$ dài $80\text{ cm}$. Đoạn $AB$ dài hơn đoạn $CD$ là
\choice
{\True $20\text{ cm}$}
{$10\text{ cm}$}
{$30\text{ cm}$}
{$80\text{ cm}$}
\loigiai{
$1\text{ m} = 100\text{ cm}$.\\
$100 - 80 = 20\text{ cm}$.\\
\textbf{Đáp án của câu hỏi là A}
}
\end{ex}

\begin{ex} %[2H11V57]
Sợi dây thứ nhất dài $2\text{ m}$, sợi dây thứ hai dài $15\text{ dm}$. Cả hai sợi dây dài tất cả là
\choice
{$30\text{ dm}$}
{$40\text{ dm}$}
{\True $35\text{ dm}$}
{$25\text{ dm}$}
\loigiai{
$2\text{ m} = 20\text{ dm}$.\\
$20 + 15 = 35\text{ dm}$.\\
\textbf{Đáp án của câu hỏi là C}
}
\end{ex}

\begin{ex} %[2H11V57]
Bạn Nam bước đi 10 bước, mỗi bước dài khoảng $4\text{ dm}$. Nam đã đi được quãng đường dài
\choice
{$30\text{ dm}$}
{\True $40\text{ dm}$ ($4\text{ m}$)}
{$50\text{ dm}$}
{$4\text{ dm}$}
\loigiai{
$4 \times 10 = 40\text{ dm} = 4\text{ m}$.\\
\textbf{Đáp án của câu hỏi là B}
}
\end{ex}

\begin{ex} %[2H11V57]
Một mảnh vườn hình chữ nhật có chiều dài $20\text{ m}$, chiều rộng $10\text{ m}$. Nửa chu vi mảnh vườn là
\choice
{\True $30\text{ m}$}
{$25\text{ m}$}
{$35\text{ m}$}
{$40\text{ m}$}
\loigiai{
$20 + 10 = 30\text{ m}$.\\
\textbf{Đáp án của câu hỏi là A}
}
\end{ex}

\begin{ex} %[2H11V57]
Bác Ba trồng 5 cây cau thẳng hàng cách đều nhau, khoảng cách giữa 2 cây liên tiếp là $2\text{ m}$. Khoảng cách từ cây đầu tiên đến cây cuối cùng là
\choice
{$10\text{ m}$}
{\True $8\text{ m}$ (có 4 khoảng cách)}
{$6\text{ m}$}
{$12\text{ m}$}
\loigiai{
5 cây có $5 - 1 = 4$ khoảng cách.\\
$2 \times 4 = 8\text{ m}$.\\
\textbf{Đáp án của câu hỏi là B}
}
\end{ex}

\begin{ex} %[2H11C57]
Một con ốc sên bò lên một cây cột cao $5\text{ m}$. Ban ngày bò lên được $3\text{ m}$, ban đêm bị tụt xuống $1\text{ m}$. Hỏi sau bao nhiêu ngày đêm ốc sên bò lên đến đỉnh cột?
\choice
{$3$ ngày đêm}
{$4$ ngày đêm}
{\True $2$ ngày đêm}
{$5$ ngày đêm}
\loigiai{
Ngày 1 lên 3m, đêm 1 tụt 1m còn 2m.\\
Ngày 2 bò tiếp 3m: $2 + 3 = 5\text{ m}$ (đạt đỉnh luôn trong ngày thứ 2).\\
\textbf{Đáp án của câu hỏi là C}
}
\end{ex}

\textbf{Phần 2. Trắc nghiệm đúng sai}
\setcounterpageref{ex}{0}

\begin{ex} %[2H11H57]
Xét tính thực tế của các đơn vị đo độ dài:
\choiceTF
{\True Quyển sách Toán 2 dài khoảng $24\text{ cm}$.}
{\True Sân trường dài khoảng vài chục mét.}
{Khoảng cách từ Hà Nội đến Đà Nẵng dài khoảng $800\text{ m}$.}
{\True Khoảng cách từ Hà Nội đến Đà Nẵng dài khoảng gần $800\text{ km}$.}
\loigiai{
\begin{itemize}
    \item Quyển sách dài $24\text{ cm}$: Đúng.
    \item Sân trường dài vài chục mét: Đúng.
    \item Khoảng cách xa phải dùng $\text{km}$ (không phải $\text{m}$). Ý 3 Sai, Ý 4 Đúng.
\end{itemize}
\begin{itemize}
    \item a) \textbf{Đúng}.
    \item b) \textbf{Đúng}.
    \item c) \textbf{Sai}.
    \item d) \textbf{Đúng}.
\end{itemize}
}
\end{ex}

\begin{ex} %[2H11V57]
Sợi dây $A$ dài $1\text{ m}$, sợi dây $B$ dài $70\text{ cm}$. Xét tính đúng sai:
\choiceTF
{\True Sợi dây $A$ dài $100\text{ cm}$.}
{\True Sợi dây $A$ dài hơn sợi dây $B$ là $30\text{ cm}$ ($100 - 70 = 30$).}
{Tổng độ dài hai sợi dây là $80\text{ cm}$.}
{\True Tổng độ dài hai sợi dây là $170\text{ cm}$ ($100 + 70 = 170$).}
\loigiai{
\begin{itemize}
    \item $1\text{ m} = 100\text{ cm}$: Đúng.
    \item Hiệu: $100 - 70 = 30\text{ cm}$: Đúng.
    \item Tổng: $170\text{ cm} \neq 80\text{ cm}$. Ý 3 Sai, Ý 4 Đúng.
\end{itemize}
\begin{itemize}
    \item a) \textbf{Đúng}.
    \item b) \textbf{Đúng}.
    \item c) \textbf{Sai}.
    \item d) \textbf{Đúng}.
\end{itemize}
}
\end{ex}

\textbf{Phần 3. Trả lời ngắn}
\setcounterpageref{ex}{0}

\begin{ex} %[2H11H57]
Đổi: $2\text{ m} = \dots\text{ cm}$.
\shortans{200}
\loigiai{
\textbf{Đáp số:} $200$\\
200 cm.
}
\end{ex}

\begin{ex} %[2H11H57]
Đổi: $3\text{ m} = \dots\text{ dm}$.
\shortans{30}
\loigiai{
\textbf{Đáp số:} $30$\\
30 dm.
}
\end{ex}

\begin{ex} %[2H11V57]
Một đoạn đường dài $45\text{ m}$, đã rải nhựa được $25\text{ m}$. Hỏi còn bao nhiêu mét đường chưa rải nhựa?
\shortans{20}
\loigiai{
\textbf{Đáp số:} $20$\\
$45 - 25 = 20\text{ m}$.
}
\end{ex}

\begin{ex} %[2H11V57]
Quãng đường từ nhà An đến trường dài $800\text{ m}$, từ trường đến nhà bà ngoại dài thêm $200\text{ m}$. Tính quãng đường từ nhà An qua trường đến nhà bà ngoại theo mét.
\shortans{1000}
\loigiai{
\textbf{Đáp số:} $1000$\\
$800 + 200 = 1000\text{ m}$.
}
\end{ex}

\begin{ex} %[2H11V57]
Có 4 đoạn dây bằng nhau, mỗi đoạn dài $5\text{ m}$. Tính tổng độ dài 4 đoạn dây theo mét.
\shortans{20}
\loigiai{
\textbf{Đáp số:} $20$\\
$5 \times 4 = 20\text{ m}$.
}
\end{ex}

\begin{ex} %[2H11C57]
Bác thợ mộc có cây gỗ dài $8\text{ m}$, bác cưa thành các đoạn dài $2\text{ m}$. Hỏi bác phải cưa bao nhiêu lần?
\shortans{3}
\loigiai{
\textbf{Đáp số:} $3$\\
Số đoạn: $8 : 2 = 4$ đoạn.\\
Số lần cưa: $4 - 1 = 3$ lần.
}
\end{ex}
